\input{../tex-inputs/latex-leseansicht-vorspann}

\section[Arthur Schnitzler an Elsa Ginsberg-Plessner, 12. 1. 1916]{L04339 Arthur Schnitzler an Elsa Ginsberg-Plessner, 12. 1. 1916}
\nopagebreak\mylabel{L04339v}
\rehead{ }\normalsize\beginnumbering\briefempfaengerindex{Plessner, Elsa@\textsc{Plessner, Elsa}!zzzSchnitzler, Arthur@\emph{von Arthur Schnitzler}!1916-01-121@{12. 1. 1916}|(be}
\toendnotes[C]{\smallbreak\pagebreak[2]}
\correspDesc{Versand  durch Arthur Schnitzler am 12. 1. 1916 in Wien
\newline{}Erhalt  durch Elsa Plessner im Zeitraum [13. 1. 1916
                  – 17. 1. 1916?] in München}\toendnotes[C]{\smallbreak}
\Standort{DLA, A:Schnitzler, HS.1985.1.829.}
\physDesc{Brief, 1 Blatt, 2 Seiten, 845 Zeichen
\newline{}Schreibmaschine
\newline{}Handschrift Schreibkraft: Bleistift (\noindent{}Ergänzungen von Satzzeichen)
\newline{}Handschrift Arthur Schnitzler: roter Buntstift, lateinische Kurrent (\noindent{}Vermerk »Ginsberg« und drei Unterstreichungen)}\toendnotes[C]{\smallbreak}
\pstart
           \raggedleft{}{\pb}12. 1. 1916.\pend
           
\pstart\center{}Sehr geehrte gnädige Frau.\pend\vspace{0.5em}
\pstart
           Eben langen \label{K_L04339-1v}\edtext{drei Couverts}{\lemma{\textnormal{\emph{drei Couverts}}}\Cendnote{\textnormal{Wie aus Plessners\pwindex{Plessner, Elsa 22.\,8.\,1875 Wien – 7.\,5.\,1932 Alicante@\textsc{Plessner, Elsa} (22.\,8.\,1875 Wien – 7.\,5.\,1932 Alicante), \emph{Schriftstellerin, Schauspielerin}|pwk} Antwortschreiben vom XXXX Auszeichnungsfehler: Dokument L03731 nicht gefunden hervorgeht, befand sich in den Couverts das nicht überlieferte Werkmanuskript von \emph{Musik}\pwindex{Plessner, Elsa 22.\,8.\,1875 Wien – 7.\,5.\,1932 Alicante@\textsc{Plessner, Elsa} (22.\,8.\,1875 Wien – 7.\,5.\,1932 Alicante), \emph{Schriftstellerin, Schauspielerin}!Musik@\strich\emph{Musik}|pwk}, im mit
                     »I« markierten Couvert lag zusätzlich ihr Anschreiben
                     (XXXX Auszeichnungsfehler: Dokument L03730 nicht gefunden).}}}\label{K_L04339-1} von Ihnen
               rekommandiert bei mir an, in denen ich wahrscheinlich nicht mit Unrecht ein Manuscript\pwindex{Plessner, Elsa 22.\,8.\,1875 Wien – 7.\,5.\,1932 Alicante@\textsc{Plessner, Elsa} (22.\,8.\,1875 Wien – 7.\,5.\,1932 Alicante), \emph{Schriftstellerin, Schauspielerin}!Musik@\strich\emph{Musik}|pwv} vermute. Sollte
               dies der Fall sein, so möchte ich Sie recht sehr bitten mich von der Lektüre entheben
               zu wollen. Aus Gründen, deren Erörterung zu weit führen würde, habe ich es endlich
               verschworen Manuscripte zu lesen, zu prüfen, zu beurteilen; und ich möchte auch Ihr
               Neues Werk\pwindex{Plessner, Elsa 22.\,8.\,1875 Wien – 7.\,5.\,1932 Alicante@\textsc{Plessner, Elsa} (22.\,8.\,1875 Wien – 7.\,5.\,1932 Alicante), \emph{Schriftstellerin, Schauspielerin}!Musik@\strich\emph{Musik}|pwv}, was es immer sein
               mag, erst kennen lernen, bis es gedruckt vorliegt. Die drei Couverts bleiben also bis
               zum Eintreffen Ihrer Antwort verschlossen bei mir liegen.\pend
           
\pstart
           Von Ziegel\pwindex{Ziegel, Erich 26.\,8.\,1876 Schwerin – 30.\,11.\,1950 München@\textsc{Ziegel, Erich} (26.\,8.\,1876 Schwerin – 30.\,11.\,1950 München), \emph{Theaterleiter, Regisseur, Schauspieler}|pw} ist bisher keinerlei Anfrage an
               mich gerichtet worden. Entschuldigen Sie bei dieser Gelegenheit, dass ich Ihnen den
               Empfang Ihres liebenswürdigen \label{K_L04339-2v}\edtext{Schreibens}{\lemma{\textnormal{\emph{Schreibens}}}\Cendnote{\textnormal{XXXX Auszeichnungsfehler: Dokument L03729 nicht gefunden.}}}\label{K_L04339-2} erst {\pb}heute mit herzlichem Dank bestätige.\pend
           
\pstart
           Mit verbindlichen Grüssen{\\[\baselineskip]}Ihr sehr ergebener\pend
           \leftskip=0em{}\selectlanguage{ngerman}\endnumbering\briefempfaengerindex{Plessner, Elsa@\textsc{Plessner, Elsa}!zzzSchnitzler, Arthur@\emph{von Arthur Schnitzler}!1916-01-121@{12. 1. 1916}|)be}\mylabel{L04339h}  \newcommand{\dateiname}{L04339}\newcommand{\titel}{Arthur Schnitzler an Elsa Ginsberg-Plessner, 12. 1. 1916}\newcommand{\editorInnen}{Herausgegeben von Jahnke, SelmaMüller, Martin Anton}\input{../tex-inputs/latex-leseansicht-abspann}
